%% =================================================================
%% Wei, Mei, Zhao, Xiao, Sun, Zhang, Guo.
%% "Nondestructively identifying the mechanical behavior of soft
%%  tissues using surface deformation with an explicit inverse
%%  approach."
%% Appl. Math. Modelling 134 (2024) 126--147.
%%
%% Reconstruction of page 17 (printed p. 142) as study notes.
%% Generated by: reproduce-basic-paper-tex → reproduce-multiple-pages-basic-tex [17, 22]
%%
%% Structure: Fig. 22 (optimization trajectory, 4 inclusions → 1)
%% + one prose paragraph (cut off at "328 iterations,")
%% + Fig. 23 (Sample 1 explicit vs implicit, no noise).
%% No equations, no tables.
%% =================================================================

\documentclass[11pt]{article}

\usepackage[margin=1in]{geometry}
\usepackage{amsmath, amssymb}
\usepackage{mathrsfs}
\usepackage{bm}
\usepackage{xcolor}
\usepackage{fancyhdr}

\pagestyle{fancy}
\fancyhf{}
\lhead{\small Y.~Mei et al.}
\rhead{\small Applied Mathematical Modelling 134 (2024) 126--147}
\cfoot{\small 142 $\to$ \arabic{page}}
\renewcommand{\headrulewidth}{0.4pt}

\setcounter{page}{1}

\begin{document}

% ---------- Figure 22 placeholder ----------
\noindent
\begin{center}
\fbox{\parbox{0.92\textwidth}{\centering\vspace{1em}
\textbf{[Figure 22 --- placeholder]}\\[0.8em]
Six-panel optimization trajectory for the \textbf{Case 3 Sample 1}
(one embedded inclusion, target SM $\approx 5$\,MPa against a
1\,MPa soft-tissue background), using the explicit inverse scheme
with 3 surface displacement datasets (no noise). The initial
guess has \textbf{four} inclusions predefined at the four corners
of the square domain; the optimization merges them into a single
central inclusion as it converges.\\[0.4em]
\begin{tabular}{@{}ll@{}}
\textbf{(a) Initial guess}    & 4 green circles at corners, SM $\in [1.00, 1.00]$\,MPa (uniform) \\
\textbf{(b) Iteration 5}      & 4 merging blobs, SM $\in [1.00, 3.20]$\,MPa \\
\textbf{(c) Iteration 10}     & 3 blobs (center two merging), SM $\in [1.00, 3.23]$\,MPa \\
\textbf{(d) Iteration 15}     & 1--2 blobs (red central), SM $\in [1.00, 4.26]$\,MPa \\
\textbf{(e) Iteration 275}    & 1 red inclusion, SM $\in [2.10, 5.32]$\,MPa \\
\textbf{(f) Iteration 328 (final)} & 1 red inclusion, SM $\in [1.00, 5.10]$\,MPa
\end{tabular}\\[0.4em]
\textbf{Takeaway}: the 4-inclusion initialization does NOT cause
the algorithm to find 4 tumors. Instead, the four components
merge into the single ground-truth tumor by iteration 15. This
is the opposite of the usual over-parameterization concern:
redundant initial components collapse to the correct solution
rather than forcing a spurious multi-component reconstruction.
The explicit parameterization has 69 total variables
(4 inclusions $\times$ 16 geometric parameters + 4 shear moduli +
1 background shear modulus).
\vspace{1em}}}
\end{center}
\vspace{0.3em}
\begin{center}
\textbf{Fig. 22.} The optimization process with four inclusions
initially utilizing the explicit inverse scheme for a solid with
inclusion embedded when 3 surface displacement datasets are
utilized (Without noise, Unit: MPa).
\end{center}

\vspace{0.8em}

% ---------- Prose paragraph (continuation from page 16) ----------
\noindent
domain. Each inclusion was represented by 16 geometric parameters
and 1 material property, which is the shear modulus. Consequently,
the total number of optimization variables reached 69, inclusive
of the shear modulus value of the background.

During the optimization process, the four predefined components
gradually merged into one single inclusion. As convergence was
reached, the shear modulus of this inclusion approached the
target values effectively, demonstrating the success of the
inverse problem-solving method. The reconstructed results, as
shown in Fig.~22, illustrated the accurate estimation of the
shear modulus for the inclusion after the convergence of the
algorithm. Indeed, the explicit inverse method demonstrated a
notably faster convergence rate compared to the implicit inverse
method. Specifically, the explicit inverse method achieved
convergence in just 328 iterations,

\vspace{0.8em}

% ---------- Figure 23 placeholder ----------
\noindent
\begin{center}
\fbox{\parbox{0.92\textwidth}{\centering\vspace{1em}
\textbf{[Figure 23 --- placeholder]}\\[0.8em]
Six-panel comparison of the \textbf{Explicit Inverse Method} (top
row) vs.\ \textbf{Implicit Inverse Method} (bottom row) on
\textbf{Case 3 Sample 1} (one embedded inclusion, target
$\approx 5$\,MPa), with \textbf{no noise}, using 3, 6, and 9
surface displacement datasets.\\[0.4em]
\textbf{Explicit (top row)}:
(a) 3 fields, SM $\in [1.00, 5.10]$\,MPa;
(b) 6 fields, SM $\in [1.00, 5.06]$\,MPa;
(c) 9 fields, SM $\in [1.00, 4.97]$\,MPa.
All three panels recover a clean single red inclusion on a blue
background, near the 5\,MPa nominal target. The colorbar max
drifts slightly downward as more fields are added
($5.10 \to 5.06 \to 4.97$).\\[0.4em]
\textbf{Implicit (bottom row)}:
(d) 3 fields, SM $\in [1.00, 4.47]$\,MPa;
(e) 6 fields, SM $\in [1.00, 4.58]$\,MPa;
(f) 9 fields, SM $\in [1.00, 4.79]$\,MPa.
\textbf{All three panels undershoot the 5\,MPa nominal target},
with the 3-field case falling short by 10.6\,\%. The max values
are non-monotone in dataset count
($4.47 \to 4.58 \to 4.79$) but monotonically approach the target
as more fields are added --- the opposite direction from the
explicit method's slight drift.\\[0.3em]
\textbf{Takeaway}: on the one-inclusion case without noise, both
methods produce visually similar reconstructions, but the
explicit method hits the target nominal from above
(slight overshoot) while the implicit method hits it from below
(persistent undershoot). More data helps the implicit method
approach nominal but doesn't quite reach it.
\vspace{1em}}}
\end{center}
\vspace{0.3em}
\begin{center}
\textbf{Fig. 23.} The reconstructed results with varying 3, 6 and
9 surface displacement datasets corresponding to Sample 1
(Without noise, Unit: MPa).
\end{center}

% [page ends here — continues on page 18]

\end{document}
